\subsection{Stiffness--solver map and loss landscape}
\label{sec:track-d}

Numerical stiffness (a system whose fastest and slowest timescales differ by
orders of magnitude) is where explicit integrators are classically expected
to struggle. We sweep damping ratio
$\mu\in\{0.1,0.5,1.0,2.0,5.0,8.0\}$ against four explicit solvers \{Euler,
Midpoint, RK4, Tsit5\} \cite{tsitouras2011tsit5} on the damped pendulum, at a
fixed step $\Delta t=0.05$. The goal is to map which $(\text{solver},\mu)$
combinations converge and which fail, and to test whether any failure is
actually a stiffness effect.

\begin{expectbox}
\expectrow{Question}{As damping ratio increases the pendulum's dynamics stiffen -- do lower-order explicit solvers genuinely fail as stiffness grows?}
\expectrow{Expected}{Higher damping ratios should push some (solver, $\mu$) combinations outside that solver's classical stability region, producing divergence.}
\expectrow{Observed}{Linear stability theory predicts zero divergences at every tested $\mu$, and the sweep confirms it. Every observed failure traces to random seed, not stiffness.}
\end{expectbox}

\subsubsection{The stability question resolves quickly}

Linearizing at equilibrium gives eigenvalues
$\lambda = \big({-\mu\pm\sqrt{\mu^2-4g/L}}\big)/2$. For every underdamped
$\mu<2\sqrt{g/L}\approx6.26$, $|\lambda|=\sqrt{g/L}\approx3.13$ is constant, so
$h|\lambda|\approx0.157$ at $\Delta t=0.05$, far inside even Euler's
stability region $h|\lambda|\le2$ \cite{dahlquist1963stability}. Even the
overdamped case ($\mu=8.0$, fast eigenvalue $\lambda_{\text{fast}}\approx-7.7$)
gives $h|\lambda_{\text{fast}}|\approx0.385$, still comfortably stable. This
predicts, and the sweep confirms, zero divergences across all 24 cells: not
one run produced an invalid gradient at any point in training. All four
solvers tested are explicit Runge-Kutta methods with a bounded stability
region. We implement no implicit solver, so this analysis cannot say whether
an implicit method would behave differently at a genuinely stiffer
configuration.

\subsubsection{What actually drives the failures: initialization}

\begin{table*}[t]
\centering
\footnotesize
\renewcommand{\arraystretch}{1.15}
\caption{Original 24-cell sweep (single seed 42), best training MSE and verdict.}
\label{tab:stiffness-single-seed}
\begin{tabularx}{\textwidth}{@{}C{1.6cm} X X X X@{}}
\toprule
\hdrrow \thh{Damping ratio $\mu$} & \thh{Euler} & \thh{Midpoint} & \thh{RK4} & \thh{Tsit5}\\
\midrule
0.1 & unstable, $1.7\times10^{-2}$ & unstable, $3.0\times10^{-2}$ & unstable, $2.2\times10^{-2}$ & unstable, $2.4\times10^{-2}$\\
0.5 & \textbf{converged}, $4.6\times10^{-3}$ & \textbf{converged}, $7.9\times10^{-3}$ & \textbf{converged}, $4.7\times10^{-3}$ & \textbf{converged}, $9.2\times10^{-3}$\\
1.0 & unstable, $1.0\times10^{-1}$ & converged, $5.4\times10^{-4}$ & converged, $5.3\times10^{-4}$ & converged, $4.0\times10^{-4}$\\
2.0 & unstable, $1.8\times10^{-2}$ & unstable, $1.8\times10^{-2}$ & unstable, $1.8\times10^{-2}$ & unstable, $1.8\times10^{-2}$\\
5.0 & converged, $6.1\times10^{-5}$ & converged, $7.0\times10^{-5}$ & converged, $6.9\times10^{-5}$ & converged, $6.8\times10^{-5}$\\
8.0 & converged, $2.1\times10^{-5}$ & converged, $2.3\times10^{-5}$ & converged, $2.6\times10^{-5}$ & converged, $2.6\times10^{-5}$\\
\bottomrule
\end{tabularx}
\end{table*}

\begin{figure}[tb]
\centering
\includegraphics[width=0.85\columnwidth]{figures/track_d/stability_heatmap_dt0.05.png}
\caption{Stability verdict heatmap, 24-cell single-seed probe sweep, $\Delta t=0.05$.}
\label{fig:stability-heatmap}
\end{figure}

Two anomalies stand out: $\mu=2.0$ fails identically across all four solvers
to two significant figures, and $\mu=1.0$ fails only for Euler. Retraining
Euler at $\mu=2.0$ with two more seeds isolates the cause (seed 42: unstable,
$1.82\times10^{-2}$; seed 1: unstable, $1.68\times10^{-2}$; seed 7:
\textbf{converged}, $7.09\times10^{-5}$): the trap depends on the random
initialization. Extending this to all four solvers $\times$ six $\mu$
$\times$ three seeds (72 cells, checkpoints saved throughout) confirms the
pattern generally. Only the two extremes ($\mu=0.1$, uniformly hard;
$\mu\in\{5.0,8.0\}$, uniformly easy) hold their verdict across every seed.
Every transitional $\mu\in\{0.5,1.0,2.0\}$ has at least one of three seeds
land in a bad local minimum. At $\mu=1.0$, two \emph{different} seeds fail
for two different reasons: seed 1 produces a universally bad initialization
shared by all four solvers (MSE $\approx9.7\text{--}9.8\times10^{-2}$, the
same signature as the $\mu=2.0$ shared trap), while seed 42 carries a second,
Euler-specific failure on top of it. Euler's apparent $\mu=1.0$ weakness in
Table~\ref{tab:stiffness-single-seed} is therefore the sum of two separate
effects, one shared and one solver-specific, compressed into a single
misleading number.

\begin{keypoint}[Consequence for the original heatmap]
The 24-cell sweep used one shared seed for all four solvers. All four solvers
failing identically at $\mu=2.0$ reflects four solvers inheriting the same
coin flip from a shared initialization, at this step size the seed$\times\mu$
interaction determines outcome far more than solver choice does.
\end{keypoint}

\subsubsection{The mechanism: a genuine spurious equilibrium}

We evaluated a checkpoint saved for Euler's trapped $\mu=2.0$ run at the true
equilibrium and at a nearby grid-search minimum. The learned field
$f_\theta$ has a large nonzero derivative at $(0,0)$ ($|f_\theta|=2.74$) and a
genuine, stable, spurious equilibrium $0.05$--$0.07$ units away
($|f_\theta|=0.0086$). This holds independently for all four solvers' own
trained models ($|f_\theta(0,0)|\in[2.65,2.88]$ across euler, midpoint, rk4,
tsit5): all four independently-trained networks converge to essentially the
same wrong fixed point.

\begin{figure*}[t]
\centering
\begin{minipage}[t]{0.37\textwidth}
\centering
\includegraphics[width=\linewidth]{figures/track_d/phase_portrait_3d_mu2.0.png}
\end{minipage}\hspace{0.04\textwidth}
\begin{minipage}[t]{0.37\textwidth}
\centering
\includegraphics[width=\linewidth]{figures/track_d/phase_portrait_3d_mu1.0.png}
\end{minipage}
\caption{Learned vector-field landscape $|f_\theta(\theta,\omega)|$, per solver's
own trained model. Left ($\mu=2.0$): every solver's true equilibrium (star)
sits partway up a slope, while a separate, lower point nearby is where each
trajectory actually settles. Right ($\mu=1.0$): Euler's star sits high on a
ridge, well apart from where its trajectory spirals down and rests; for the
other three solvers, star and resting point coincide at the valley floor.}
\label{fig:phase-portraits}
\end{figure*}

The full predicted trajectories tell the same story over time rather than in
a single snapshot. Figure~\ref{fig:mu2-collage} shows every solver's
predicted angle, angular velocity, and error at $\mu=2.0$ settling onto the
same wrong, generic trajectory; Figure~\ref{fig:mu1-collage} shows the
$\mu=1.0$ case, where only Euler's predicted angle visibly departs from the
true one.

\begin{figure*}[t]
\centering
\includegraphics[width=\textwidth]{figures/track_d/mu2.0_trajectory_collage.png}
\caption{$\mu=2.0$, all four solvers: predicted angle, angular velocity, and
error over time, alongside each solver's phase portrait.}
\label{fig:mu2-collage}
\end{figure*}

\begin{figure*}[t]
\centering
\includegraphics[width=\textwidth]{figures/track_d/mu1.0_trajectory_collage.png}
\caption{$\mu=1.0$, all four solvers: the same panel layout, showing Euler's
trajectory diverging from the true angle while the other three track it
closely.}
\label{fig:mu1-collage}
\end{figure*}

\subsubsection{Weight-space loss landscape: an independent measurement}

The phase-space plots above describe the learned dynamics. A separate
question is how forgiving the \emph{training loss} is around each solution:
does a small nudge to the weights leave the fit essentially unchanged, or
make it collapse? We measured this by stepping each trained model's weights
along two random directions and re-evaluating training error on a grid,
following the visualization technique of Li et al.\ \cite{li2018losslandscape}
adapted to KAN-ODE's 360-parameter weight space. Every cell traces the same
qualitative shape (a dip at the trained point that rises moving outward),
and we summarize each by its ``rise'': the number of orders of magnitude the
log-loss climbs by the edge of the grid.

\begin{table}[tb]
\centering
\small
\renewcommand{\arraystretch}{1.1}
\caption{Loss-landscape rise, single direction-pair (seed 0).}
\label{tab:rise-seed0}
\begin{tabularx}{\columnwidth}{@{}C{0.9cm} Y Y Y Y@{}}
\toprule
\hdrrow \thh{$\mu$} & \thh{euler} & \thh{midpoint} & \thh{rk4} & \thh{tsit5}\\
\midrule
1.0 & \textbf{5.53} & 9.03 & 8.99 & 9.12\\
2.0 & \textbf{4.23} & \textbf{4.24} & \textbf{4.22} & \textbf{4.20}\\
5.0 & 8.75 & 8.75 & 8.76 & 8.76\\
8.0 & 9.14 & 9.15 & 9.12 & 9.12\\
\bottomrule
\end{tabularx}
\end{table}

The converged models (rise $\approx9$) sit at the bottom of a narrow, steep
basin: stepping even a small distance from the trained weights in either
random direction drives the log-loss up by nine orders of magnitude, so the
surrounding terrain looks like a sharp valley with the solution pinned at its
floor. The trapped cells (bold, rise $\approx4$--$5.5$) sit in a basin roughly
half as steep. The terrain around them is closer to a broad, gently sloped
plain than a valley, and the same size step barely raises the loss by
comparison. This is the opposite of the usual intuition that a flat minimum
generalizes better. Here, the flat region is where training settled for a
generic, low-effort answer, and the steep, narrow basin is where it actually
learned the pendulum's dynamics.

Figure~\ref{fig:landscape-mu2-euler} renders this directly as a surface. The
terrain around Euler's trapped $\mu=2.0$ solution is a broad, gently rolling
plain with no sharply defined bottom. This contrasts clearly with
Figure~\ref{fig:landscape-mu1-tsit5}'s converged $\mu=1.0$ solution, which
sits at the bottom of a narrow, sharply walled valley; it is the same
contrast the rise numbers above report as a factor of roughly two, made
visible. Figure~\ref{fig:landscape-mu1-euler} shows the trapped $\mu=1.0$
case under the same solver whose $\mu=2.0$ terrain is shown alongside it,
confirming the flat terrain is a property of the trapped solution rather than
of any one damping value.

\begin{figure*}[t]
\centering
\includegraphics[width=0.75\textwidth]{figures/track_d/mu2.0_euler_direction_comparison.png}
\caption{Euler, $\mu=2.0$ (trapped): the training-loss surface around the
trained weights, viewed under three independent random direction-pairs. The
basin near the trained point (the marked star) is broad and shallow in all
three views; the third panel's steep spike on the far side is the outlier
direction discussed below, not the basin itself.}
\label{fig:landscape-mu2-euler}
\end{figure*}

\begin{figure*}[t]
\centering
\includegraphics[width=0.75\textwidth]{figures/track_d/mu1.0_euler_direction_comparison.png}
\caption{Euler, $\mu=1.0$ (also trapped): the same broad, shallow basin shape
recurs at a different damping value under the same solver.}
\label{fig:landscape-mu1-euler}
\end{figure*}

\begin{figure*}[t]
\centering
\includegraphics[width=0.75\textwidth]{figures/track_d/mu1.0_tsit5_direction_comparison.png}
\caption{Tsit5, $\mu=1.0$ (converged): the basin around the trained point is
visibly narrower and drops noticeably lower, in every one of the three
direction-pairs, than either trapped case above.}
\label{fig:landscape-mu1-tsit5}
\end{figure*}

\paragraph{Robustness check across three direction-pairs.} Repeating the
measurement with two more independent random direction pairs per cell (72
measurements total) found one direction-pair acting as a shared outlier
across nearly every cell. Because every model shares the same architecture,
one particular raw random draw runs toward a steeper wall in all 24 cells at
once, inflating that direction's rise by roughly $2$--$3\times$ everywhere.
The size of the measured effect depends on how many direction choices are
averaged. The single-seed $\approx2\times$ gap between trapped and converged
cells narrows to roughly 20--30\% once all three direction-pairs are
averaged, since the outlier direction inflates the trapped cells' mean more
than the converged cells' mean. The ranking between cells holds regardless;
only the apparent size of the gap changes. The multi-seed average is the
figure to trust; the single-seed table above illustrates the shape of the
effect at a scale that is easier to see.

\begin{figure*}[t]
\centering
\includegraphics[width=0.65\textwidth]{figures/track_d/rise_comparison_bars.png}
\caption{Mean rise per (solver, $\mu$), 3-seed average with seed-to-seed range
as error bars, colored by verdict.}
\label{fig:rise-bars}
\end{figure*}

The cells trapped by the spurious-equilibrium mechanism ($\mu=2.0$ across
all solvers, and Euler's own $\mu=1.0$) sit around 20--30\% lower on this
chart than converged cells: a modest but measured gap rather than an
order-of-magnitude separation. $\mu=0.1$ carries the ``unstable'' label for
an unrelated reason: its trajectory never settles within the training window
at all, and its bars sit in the same range as the converged
$\mu\in\{5.0,8.0\}$ cells, occasionally above them. The flatness effect
belongs specifically to the spurious-equilibrium trap cases, and does not
generalize to every cell the pipeline labels ``unstable.''

\begin{figure*}[t]
\centering
\begin{minipage}[t]{0.36\textwidth}
\centering
\includegraphics[width=\linewidth]{figures/track_d/seed_converged_fraction_heatmap.png}
\end{minipage}\hspace{0.04\textwidth}
\begin{minipage}[t]{0.36\textwidth}
\centering
\includegraphics[width=\linewidth]{figures/track_d/seed_median_mse_heatmap.png}
\end{minipage}
\caption{Converged fraction (left) and median training MSE (right) over 3
seeds, all 4 solvers $\times$ 6 $\mu$ (72 cells). The two agree everywhere:
low converged-fraction cells are exactly the cells with elevated median
error.}
\label{fig:seed-converged-fraction}
\end{figure*}

\begin{keypoint}[Verdict]
Every divergence-free prediction from linear stability theory held exactly,
so nothing in this sweep is a classical numerical-stiffness failure. The
mechanism at work is a seed-dependent local minimum in a 360-parameter
optimization landscape, producing a genuine, solver-independent spurious
equilibrium in the learned dynamics. Three independent lines of evidence
converge on this: per-solver trained-weight evaluation, phase-space
vector-field visualization, and an independently-measured, direction-robust
weight-space loss landscape.
\end{keypoint}
